\section{Related Work}

\yuji{In related work, we need to clarify several points: 1. why we use parameter-efficient fine-tuning (minimal disturbance of model utility and exploration into ripple effects under small updates (or clarify it in the setting section?) 2. the difference between expected ripple effect and unexpected ripple effect.)}

\subsection{Knowledge Injection in LLMs}
Knowledge updating has become a key approach to enhance LLMs' factuality, with methods spanning fine-tuning \cite{zhu2020modifying}, prompt engineering \cite{wei2022chain}, and knowledge graph (KG) integration \cite{pan2024roadmap, yasunaga2021qagnn}. 


\citet{bosselut2019comet} fine-tuned LLMs with KG triples to improve question-answering accuracy, while \citet{yasunaga2021qagnn} integrated KG embeddings into LLM layers to reduce hallucinations. 

\yuji{Revise this paragraph to emphasize that previous work focuses more on expected ripple effects and explores more about how similarity or logic affects ripple effects}
However, most studies only evaluate post-injection performance via task accuracy (e.g., answer correctness) and rarely investigate how \emph{false}\yuji{not necessarily false, common knowledge updating also causes it.} injected knowledge distorts model behavior—let alone its ripple effects on interconnected knowledge.


\subsection{Calibration of LLM Confidence}
\yuji{We probably need to delete this section to avoid distraction to confidence}
Token-level confidence scores (e.g., softmax probabilities) are widely used to assess LLM reliability, but research has revealed widespread calibration issues. 

% "Showed that LLMs are overconfident..."
\citet{kadavath2022language} showed that LLMs are overconfident in out-of-distribution scenarios, while \citet{guo2017calibration} proposed metrics like Expected Calibration Error (ECE) to quantify mismatches between confidence and accuracy. 

Recent work has focused on improving calibration: \citet{kuhn2023semantic} explored linguistic invariances to estimate uncertainty, and other works have optimized training objectives for better confidence alignment. Yet these studies rarely link confidence miscalibration to \emph{knowledge injection}—they do not explore how false injected knowledge induces ``false confidence'' or invalidate confidence as an evaluation metric for post-injection LLMs.

\subsection{LLMs and Knowledge Graph Interaction}
\yuji{This section is irrelevant to our focus. We can merge the related work with how we construct our knowledge topology that reflects real-world distribution.}
To model LLMs' internal knowledge structure, recent work has explored connections between LLMs and KGs. 

% "LLMs implicitly encode KG-like relational knowledge" -> 这必须是 Knowledge Neurons 或 ROME
\citet{dai2022kn} showed that LLMs implicitly encode KG-like relational knowledge, while \citet{cohen2023ripple} analyzed how factual errors spread across LLM-generated text (Ripple Effects). 

However, existing studies either: (1) Treat LLMs as KG extractors (not receivers of injected knowledge) \cite{bosselut2019comet}, or (2) Only measure error propagation in unstructured text (not structured KG nodes) \cite{cohen2023ripple, yao2023editing}. No work has systematically studied how false injected knowledge ripples through KG node connections (e.g., 1-hop vs. 2-hop neighbors) or linked this propagation to token-level confidence changes—two core focuses of our research.